% Options for packages loaded elsewhere
\PassOptionsToPackage{unicode}{hyperref}
\PassOptionsToPackage{hyphens}{url}
%
\documentclass[
]{article}
\usepackage{amsmath,amssymb}
\usepackage{iftex}
\ifPDFTeX
  \usepackage[T1]{fontenc}
  \usepackage[utf8]{inputenc}
  \usepackage{textcomp} % provide euro and other symbols
\else % if luatex or xetex
  \usepackage{unicode-math} % this also loads fontspec
  \defaultfontfeatures{Scale=MatchLowercase}
  \defaultfontfeatures[\rmfamily]{Ligatures=TeX,Scale=1}
\fi
\usepackage{lmodern}
\ifPDFTeX\else
  % xetex/luatex font selection
\fi
% Use upquote if available, for straight quotes in verbatim environments
\IfFileExists{upquote.sty}{\usepackage{upquote}}{}
\IfFileExists{microtype.sty}{% use microtype if available
  \usepackage[]{microtype}
  \UseMicrotypeSet[protrusion]{basicmath} % disable protrusion for tt fonts
}{}
\makeatletter
\@ifundefined{KOMAClassName}{% if non-KOMA class
  \IfFileExists{parskip.sty}{%
    \usepackage{parskip}
  }{% else
    \setlength{\parindent}{0pt}
    \setlength{\parskip}{6pt plus 2pt minus 1pt}}
}{% if KOMA class
  \KOMAoptions{parskip=half}}
\makeatother
\usepackage{xcolor}
\usepackage{color}
\usepackage{fancyvrb}
\newcommand{\VerbBar}{|}
\newcommand{\VERB}{\Verb[commandchars=\\\{\}]}
\DefineVerbatimEnvironment{Highlighting}{Verbatim}{commandchars=\\\{\}}
% Add ',fontsize=\small' for more characters per line
\newenvironment{Shaded}{}{}
\newcommand{\AlertTok}[1]{\textcolor[rgb]{1.00,0.00,0.00}{\textbf{#1}}}
\newcommand{\AnnotationTok}[1]{\textcolor[rgb]{0.38,0.63,0.69}{\textbf{\textit{#1}}}}
\newcommand{\AttributeTok}[1]{\textcolor[rgb]{0.49,0.56,0.16}{#1}}
\newcommand{\BaseNTok}[1]{\textcolor[rgb]{0.25,0.63,0.44}{#1}}
\newcommand{\BuiltInTok}[1]{\textcolor[rgb]{0.00,0.50,0.00}{#1}}
\newcommand{\CharTok}[1]{\textcolor[rgb]{0.25,0.44,0.63}{#1}}
\newcommand{\CommentTok}[1]{\textcolor[rgb]{0.38,0.63,0.69}{\textit{#1}}}
\newcommand{\CommentVarTok}[1]{\textcolor[rgb]{0.38,0.63,0.69}{\textbf{\textit{#1}}}}
\newcommand{\ConstantTok}[1]{\textcolor[rgb]{0.53,0.00,0.00}{#1}}
\newcommand{\ControlFlowTok}[1]{\textcolor[rgb]{0.00,0.44,0.13}{\textbf{#1}}}
\newcommand{\DataTypeTok}[1]{\textcolor[rgb]{0.56,0.13,0.00}{#1}}
\newcommand{\DecValTok}[1]{\textcolor[rgb]{0.25,0.63,0.44}{#1}}
\newcommand{\DocumentationTok}[1]{\textcolor[rgb]{0.73,0.13,0.13}{\textit{#1}}}
\newcommand{\ErrorTok}[1]{\textcolor[rgb]{1.00,0.00,0.00}{\textbf{#1}}}
\newcommand{\ExtensionTok}[1]{#1}
\newcommand{\FloatTok}[1]{\textcolor[rgb]{0.25,0.63,0.44}{#1}}
\newcommand{\FunctionTok}[1]{\textcolor[rgb]{0.02,0.16,0.49}{#1}}
\newcommand{\ImportTok}[1]{\textcolor[rgb]{0.00,0.50,0.00}{\textbf{#1}}}
\newcommand{\InformationTok}[1]{\textcolor[rgb]{0.38,0.63,0.69}{\textbf{\textit{#1}}}}
\newcommand{\KeywordTok}[1]{\textcolor[rgb]{0.00,0.44,0.13}{\textbf{#1}}}
\newcommand{\NormalTok}[1]{#1}
\newcommand{\OperatorTok}[1]{\textcolor[rgb]{0.40,0.40,0.40}{#1}}
\newcommand{\OtherTok}[1]{\textcolor[rgb]{0.00,0.44,0.13}{#1}}
\newcommand{\PreprocessorTok}[1]{\textcolor[rgb]{0.74,0.48,0.00}{#1}}
\newcommand{\RegionMarkerTok}[1]{#1}
\newcommand{\SpecialCharTok}[1]{\textcolor[rgb]{0.25,0.44,0.63}{#1}}
\newcommand{\SpecialStringTok}[1]{\textcolor[rgb]{0.73,0.40,0.53}{#1}}
\newcommand{\StringTok}[1]{\textcolor[rgb]{0.25,0.44,0.63}{#1}}
\newcommand{\VariableTok}[1]{\textcolor[rgb]{0.10,0.09,0.49}{#1}}
\newcommand{\VerbatimStringTok}[1]{\textcolor[rgb]{0.25,0.44,0.63}{#1}}
\newcommand{\WarningTok}[1]{\textcolor[rgb]{0.38,0.63,0.69}{\textbf{\textit{#1}}}}
\setlength{\emergencystretch}{3em} % prevent overfull lines
\providecommand{\tightlist}{%
  \setlength{\itemsep}{0pt}\setlength{\parskip}{0pt}}
\setcounter{secnumdepth}{-\maxdimen} % remove section numbering
\ifLuaTeX
  \usepackage{selnolig}  % disable illegal ligatures
\fi
\usepackage{bookmark}
\IfFileExists{xurl.sty}{\usepackage{xurl}}{} % add URL line breaks if available
\urlstyle{same}
\hypersetup{
  hidelinks,
  pdfcreator={LaTeX via pandoc}}

\author{}
\date{}

\begin{document}

\section{Introdução ao Networking}\label{introduuxe7uxe3o-ao-networking}

\subsubsection{Conteúdo}\label{conteuxfado}

\begin{itemize}
\tightlist
\item
  \hyperref[introduuxe7uxe3o-ao-networking]{Introdução ao Networking}

  \begin{itemize}
  \tightlist
  \item
    \hyperref[aplicauxe7uxe3o-de-rede-simples]{Aplicação de rede
    simples}
  \item
    \hyperref[conectando-de-fora-do-cluster]{Conectando de fora do
    cluster}
  \item
    \hyperref[configurauxe7uxe3o-inicial]{Configuração inicial}

    \begin{itemize}
    \tightlist
    \item
      \hyperref[limitauxe7uxf5es-do-ambiente-local]{Limitações do
      ambiente local}
    \end{itemize}
  \item
    \hyperref[o-que-uxe9-um-serviuxe7o]{O que é um serviço?}
  \item
    \hyperref[sobre-o-roteamento-do-kubernetes]{Sobre o roteamento do
    Kubernetes}
  \item
    \hyperref[o-que-uxe9-um-ingress]{O que é um Ingress?}
  \item
    \hyperref[api-de-gateway]{API de gateway}
  \end{itemize}
\end{itemize}

Agora de volta ao desenvolvimento! Reiniciar e acompanhar os logs foi um
deleite. Em seguida, abriremos um endpoint para o aplicativo e o
acessaremos via HTTP.

\subsubsection{Aplicação de rede
simples}\label{aplicauxe7uxe3o-de-rede-simples}

Vamos desenvolver nosso aplicativo para que ele tenha um servidor HTTP
respondendo com dois hashes: um hash que é armazenado até que o processo
seja encerrado e um hash que é específico da solicitação. O corpo da
resposta pode ser algo como ``Aplicativo abc123. Solicitação 94k9m2''.
Escolha qualquer porta para ouvir.

Um
\href{https://github.com/kubernetes-hy/material-example/tree/master/app2}{exemplo}
já foi preparado. Por padrão, ele escutará na porta 3000.

\begin{Shaded}
\begin{Highlighting}[]
\ExtensionTok{$}\NormalTok{ kubectl apply }\AttributeTok{{-}f}\NormalTok{ https://raw.githubusercontent.com/kubernetes{-}hy/material{-}example/master/app2/manifests/deployment.yaml}
\ExtensionTok{deployment.apps/hashresponse{-}dep}\NormalTok{ created}
\end{Highlighting}
\end{Shaded}

\subsubsection{Conectando de fora do
cluster}\label{conectando-de-fora-do-cluster}

Podemos confirmar que o hashresponse-dep está funcionando com o comando
\texttt{port-forward}. Vamos ver primeiro o nome do pod e depois fazer o
port-forward:

\begin{Shaded}
\begin{Highlighting}[]
\ExtensionTok{$}\NormalTok{ kubectl get po}
\ExtensionTok{NAME}\NormalTok{                                READY   STATUS    RESTARTS   AGE}
\ExtensionTok{hashgenerator{-}dep{-}5cbbf97d5{-}z2ct9}\NormalTok{   1/1     Running   0          20h}
\ExtensionTok{hashresponse{-}dep{-}57bcc888d7{-}dj5vk}\NormalTok{   1/1     Running   0          19h}

\ExtensionTok{$}\NormalTok{ kubectl port{-}forward hashresponse{-}dep{-}57bcc888d7{-}dj5vk 3003:3000}
\ExtensionTok{Forwarding}\NormalTok{ from 127.0.0.1:3003 }\AttributeTok{{-}}\OperatorTok{\textgreater{}}\NormalTok{ 3000}
\ExtensionTok{Forwarding}\NormalTok{ from }\PreprocessorTok{[}\SpecialStringTok{::1}\PreprocessorTok{]}\NormalTok{:3003 }\AttributeTok{{-}}\OperatorTok{\textgreater{}}\NormalTok{ 3000}
\end{Highlighting}
\end{Shaded}

Agora podemos ver a resposta de http://localhost:3003 e confirmar que
está funcionando conforme o esperado.

Como vemos, \texttt{port-forward} é um comando usado no Kubernetes para
encaminhar uma porta local para um pod. Isso permite que você acesse
facilmente os aplicativos em execução dentro do seu cluster Kubernetes a
partir da sua máquina local. O encaminhamento de porta não é destinado
ao uso em produção, mas é muito útil para fins de depuração e
desenvolvimento.

\begin{quote}
\textbf{Exercício 1.5 --- O projeto, etapa 3}

\textbf{Instruções:} Faça o projeto responder algo a uma solicitação GET
enviada para a URL do projeto. Uma página HTML simples é boa, ou você
pode implantar algo mais complexo, como um aplicativo de página única.

Veja aqui (abre em uma nova aba) como você pode definir variáveis de
ambiente para contêineres.

Use \texttt{kubectl\ port-forward} para confirmar se o projeto está
acessível e funciona no cluster usando um navegador para acessar o
projeto.
\end{quote}

Conectar-se a um aplicativo dentro de um contêiner Docker é fácil. Se
você tiver um aplicativo escutando a porta 3000, você pode simplesmente
usar o \texttt{docker\ run} com a flag \texttt{-p\ 3003:3000}. Outra
opção é usar a declaração de portas do Docker Compose. Infelizmente, o
Kubernetes não é tão simples. Usaremos um recurso de serviço, um recurso
de entrada (Ingress) ou a solução mais recente, a API Gateway.

\subsubsection{Configuração inicial}\label{configurauxe7uxe3o-inicial}

Como estamos executando nosso cluster dentro do Docker com o k3d,
teremos que fazer alguns preparativos.

Abrir uma rota de fora do cluster para o pod não será suficiente se não
tivermos meios de acessar o cluster k3d dentro dos contêineres do
Docker! Vamos ver como é a configuração padrão:

\begin{Shaded}
\begin{Highlighting}[]
\ExtensionTok{$}\NormalTok{ docker ps}
\ExtensionTok{CONTAINER}\NormalTok{ ID    IMAGE                      COMMAND                  CREATED             STATUS              PORTS                             NAMES}
\ExtensionTok{b60f6c246ebb}\NormalTok{    ghcr.io/k3d{-}io/k3d{-}proxy:5 }\StringTok{"/bin/sh {-}c nginx{-}pr…"}\NormalTok{   2 hours ago         Up 2 hours          80/tcp, 0.0.0.0:58264{-}}\OperatorTok{\textgreater{}}\NormalTok{6443/tcp   k3d{-}k3s{-}default{-}serverlb}
\ExtensionTok{553041f96fc6}\NormalTok{    rancher/k3s:latest         }\StringTok{"/bin/k3s agent"}\NormalTok{         2 hours ago         Up 2 hours                                            k3d{-}k3s{-}default{-}agent{-}1}
\ExtensionTok{aebd23c2ef99}\NormalTok{    rancher/k3s:latest         }\StringTok{"/bin/k3s agent"}\NormalTok{         2 hours ago         Up 2 hours                                            k3d{-}k3s{-}default{-}agent{-}0}
\ExtensionTok{a34e49184d37}\NormalTok{    rancher/k3s:latest         }\StringTok{"/bin/k3s server {-}{-}t…"}\NormalTok{   2 hours ago         Up 2 hours                                            k3d{-}k3s{-}default{-}server{-}0}
\end{Highlighting}
\end{Shaded}

Vemos que o k3d preparou utilmente a porta 58264 para acessar a API do
Kubernetes por meio do balanceador de carga, que faz proxy das
solicitações para a porta 6443 no nó do servidor. Além disso, a porta 80
está aberta no balanceador de carga. Todas as solicitações ao
balanceador de carga serão enviadas por proxy para o nó servidor do
cluster.

Como vemos, nosso cluster possui atualmente um nó servidor e dois nós
agentes. Para fins de teste, também desejaremos uma porta individual
aberta para um único nó de agente. Vamos excluir nosso cluster antigo e
criar um novo com portas específicas abertas para facilitar o acesso e
os testes.

A documentação do K3D (abre em uma nova aba) nos diz como as portas são
abertas: abriremos o local 8081 para 80 em
\texttt{k3d-k3s-default-serverlb} e o local 8082 para 30080 em
\texttt{k3d-k3s-default-agent-0}. O 30080 é escolhido quase
completamente aleatoriamente, mas precisa ter um valor entre 30000-32767
para a próxima etapa:

\begin{Shaded}
\begin{Highlighting}[]
\ExtensionTok{$}\NormalTok{ k3d cluster delete}
\VariableTok{INFO}\OperatorTok{[}\DecValTok{0000}\OperatorTok{]} \ExtensionTok{Deleting}\NormalTok{ cluster }\StringTok{\textquotesingle{}k3s{-}default\textquotesingle{}}
\ExtensionTok{...}
\VariableTok{INFO}\OperatorTok{[}\DecValTok{0002}\OperatorTok{]} \ExtensionTok{Successfully}\NormalTok{ deleted cluster k3s{-}default!}

\ExtensionTok{$}\NormalTok{ k3d cluster create }\AttributeTok{{-}{-}port}\NormalTok{ 8082:30080@agent:0 }\AttributeTok{{-}p}\NormalTok{ 8081:80@loadbalancer }\AttributeTok{{-}{-}agents}\NormalTok{ 2}
\VariableTok{INFO}\OperatorTok{[}\DecValTok{0000}\OperatorTok{]} \ExtensionTok{Created}\NormalTok{ network }\StringTok{\textquotesingle{}k3d{-}k3s{-}default\textquotesingle{}}
\ExtensionTok{...}
\VariableTok{INFO}\OperatorTok{[}\DecValTok{0021}\OperatorTok{]} \ExtensionTok{Cluster} \StringTok{\textquotesingle{}k3s{-}default\textquotesingle{}}\NormalTok{ created successfully!}
\VariableTok{INFO}\OperatorTok{[}\DecValTok{0021}\OperatorTok{]} \ExtensionTok{You}\NormalTok{ can now use it like this:}
\ExtensionTok{kubectl}\NormalTok{ cluster{-}info}

\ExtensionTok{$}\NormalTok{ kubectl apply }\AttributeTok{{-}f}\NormalTok{ https://raw.githubusercontent.com/kubernetes{-}hy/material{-}example/master/app2/manifests/deployment.yaml}
\ExtensionTok{deployment.apps/hashresponse{-}dep}\NormalTok{ created}
\end{Highlighting}
\end{Shaded}

Acima, o \texttt{agent:0} e o \texttt{loadbalancer} são baseados na
documentação do k3d e na leitura do código-fonte
\href{https://github.com/k3d-io/k3d/blob/11cc7979228f304025d61254eb0c0cb2745b9444/cmd/util/filter.go\#L119}{daqui}
e
\href{https://github.com/k3d-io/k3d/blob/main/pkg/types/types.go\#L65}{aqui}.

Agora temos acesso através da porta 8081 ao nosso nó servidor (na
verdade, todos os nós) e 8082 a um dos nossos nós agentes, porta 30080.
Eles serão usados para mostrar diferentes métodos de comunicação com os
servidores.

\paragraph{Limitações do ambiente
local}\label{limitauxe7uxf5es-do-ambiente-local}

A configuração não é perfeita, teremos uma quantidade limitada de portas
disponíveis no futuro. Isso será suficiente para nossos casos de uso.

Seu sistema operacional pode suportar o uso da rede host, portanto,
nenhuma porta precisa ser aberta. No entanto, não há experiência
relatada com essa abordagem neste material.

\subsubsection{O que é um serviço?}\label{o-que-uxe9-um-serviuxe7o}

Os recursos de implantação (Deployment) gerenciam o processo de
implantação para nós, cuidando da criação dos pods. Dado que os pods do
aplicativo são efêmeros e podem ser criados ou encerrados a qualquer
momento, a comunicação com o aplicativo não deve depender da presença de
nenhum pod específico.

Os recursos de serviço (Service) são essenciais para gerenciar a
acessibilidade do aplicativo, garantindo que ele possa ser alcançado por
conexões originadas tanto de fora do cluster quanto de dentro. Esses
recursos lidam com o roteamento e o balanceamento de carga necessários
para manter uma comunicação perfeita com o aplicativo, independentemente
da natureza dinâmica da criação e do encerramento do pod.

Após a recriação do cluster, devemos criar novamente também nosso
aplicativo de exemplo:

\begin{Shaded}
\begin{Highlighting}[]
\ExtensionTok{kubectl}\NormalTok{ apply }\AttributeTok{{-}f}\NormalTok{ https://raw.githubusercontent.com/kubernetes{-}hy/material{-}example/master/app2/manifests/deployment.yaml}
\end{Highlighting}
\end{Shaded}

O \texttt{deployment.yaml} tem a seguinte aparência:

\begin{Shaded}
\begin{Highlighting}[]
\FunctionTok{apiVersion}\KeywordTok{:}\AttributeTok{ apps/v1}
\FunctionTok{kind}\KeywordTok{:}\AttributeTok{ Deployment}
\FunctionTok{metadata}\KeywordTok{:}
\AttributeTok{  }\FunctionTok{name}\KeywordTok{:}\AttributeTok{ hashresponse{-}dep}
\FunctionTok{spec}\KeywordTok{:}
\AttributeTok{  }\FunctionTok{replicas}\KeywordTok{:}\AttributeTok{ }\DecValTok{1}
\AttributeTok{  }\FunctionTok{selector}\KeywordTok{:}
\AttributeTok{    }\FunctionTok{matchLabels}\KeywordTok{:}
\AttributeTok{      }\FunctionTok{app}\KeywordTok{:}\AttributeTok{ hashresponse}
\AttributeTok{  }\FunctionTok{template}\KeywordTok{:}
\AttributeTok{    }\FunctionTok{metadata}\KeywordTok{:}
\AttributeTok{      }\FunctionTok{labels}\KeywordTok{:}
\AttributeTok{        }\FunctionTok{app}\KeywordTok{:}\AttributeTok{ hashresponse}
\AttributeTok{    }\FunctionTok{spec}\KeywordTok{:}
\AttributeTok{      }\FunctionTok{containers}\KeywordTok{:}
\AttributeTok{        }\KeywordTok{{-}}\AttributeTok{ }\FunctionTok{name}\KeywordTok{:}\AttributeTok{ hashresponse}
\AttributeTok{          }\FunctionTok{image}\KeywordTok{:}\AttributeTok{ jakousa/dwk{-}app2:b7fc18de2376da80ff0cfc72cf581a9f94d10e64}
\end{Highlighting}
\end{Shaded}

Vamos agora criar um arquivo \texttt{service.yaml} na pasta
\texttt{manifests}. Precisamos do serviço para fazer as seguintes
coisas:

\begin{itemize}
\tightlist
\item
  Declarar que queremos um Serviço
\item
  Declarar qual porta ouvir
\item
  Declarar o requerimento (seletor) para onde o pedido deverá ser
  dirigido
\item
  Declarar a porta para onde o pedido deve ser dirigido
\end{itemize}

O arquivo \texttt{service.yaml} resultante é assim:

\begin{Shaded}
\begin{Highlighting}[]
\FunctionTok{apiVersion}\KeywordTok{:}\AttributeTok{ v1}
\FunctionTok{kind}\KeywordTok{:}\AttributeTok{ Service}
\FunctionTok{metadata}\KeywordTok{:}
\AttributeTok{  }\FunctionTok{name}\KeywordTok{:}\AttributeTok{ hashresponse{-}svc}
\FunctionTok{spec}\KeywordTok{:}
\AttributeTok{  }\FunctionTok{type}\KeywordTok{:}\AttributeTok{ NodePort}
\AttributeTok{  }\FunctionTok{selector}\KeywordTok{:}
\AttributeTok{    }\FunctionTok{app}\KeywordTok{:}\AttributeTok{ hashresponse}\CommentTok{ \# This is the app as declared in the deployment.}
\AttributeTok{  }\FunctionTok{ports}\KeywordTok{:}
\AttributeTok{    }\KeywordTok{{-}}\AttributeTok{ }\FunctionTok{name}\KeywordTok{:}\AttributeTok{ http}
\AttributeTok{      }\FunctionTok{nodePort}\KeywordTok{:}\AttributeTok{ }\DecValTok{30080}\CommentTok{ \# This is the port that is available outside. Value for nodePort can be between 30000{-}32767}
\AttributeTok{      }\FunctionTok{protocol}\KeywordTok{:}\AttributeTok{ TCP}
\AttributeTok{      }\FunctionTok{port}\KeywordTok{:}\AttributeTok{ }\DecValTok{1234}\CommentTok{ \# This is a port that is available to the cluster, in this case it can be \textasciitilde{} anything}
\AttributeTok{      }\FunctionTok{targetPort}\KeywordTok{:}\AttributeTok{ }\DecValTok{3000}\CommentTok{ \# This is the target port}
\end{Highlighting}
\end{Shaded}

\begin{Shaded}
\begin{Highlighting}[]
\ExtensionTok{$}\NormalTok{ kubectl apply }\AttributeTok{{-}f}\NormalTok{ manifests/service.yaml}
\ExtensionTok{service/hashresponse{-}svc}\NormalTok{ created}
\end{Highlighting}
\end{Shaded}

Como publicamos 8082 como 30080, podemos acessá-lo agora via
http://localhost:8082 .

Agora definimos um serviço com \texttt{type:\ NodePort}. Portas Node
(abre em uma nova aba) são simplesmente portas abertas pelo Kubernetes
para todos os nós, e o serviço manipulará solicitações nessa porta.
NodePorts não são flexíveis e exigem que você atribua uma porta
diferente para cada aplicativo. Como tal, os NodePorts não são usados em
produção, mas são úteis de conhecer.

O que gostaríamos de usar em vez do NodePort seria um serviço do tipo
LoadBalancer, mas isso ``só'' funciona com provedores de nuvem, pois
configura um balanceador de carga possivelmente caro para ele.
Conheceremos esse tipo no Capítulo 4.

Vamos analisar mais uma vez a configuração do aplicativo de exemplo:

\textbf{Implantação (hashresponse-dep)} - Executa 1 réplica do contêiner
\texttt{jakousa/dwk-app2} - Pod rotulado com \texttt{app:\ hashresponse}
- O contêiner escuta na porta 3000

\textbf{Serviço (hashresponse-svc)} - Tipo: NodePort --- expõe o
aplicativo externamente - Seletor: \texttt{app:\ hashresponse} ---
roteia o tráfego para pods correspondentes - Configuração da porta: -
\texttt{targetPort:\ 3000} --- a porta em que o contêiner escuta -
\texttt{port:\ 1234} --- a porta exposta pelo Serviço para acesso
interno ao cluster - \texttt{nodePort:\ 30080} --- a porta exposta em
cada nó do cluster para acesso externo

Existem duas maneiras de acessar o aplicativo:

\begin{itemize}
\tightlist
\item
  \textbf{De dentro do cluster:} \texttt{hashresponse-svc:1234}
\item
  \textbf{Da sua máquina local:} \texttt{localhost:8082} (o k3d mapeia a
  porta do host 8082 → NodePort 30080)
\end{itemize}

O tráfego flui da seguinte forma quando o aplicativo é acessado a partir
da sua máquina local:

\begin{verbatim}
localhost:8082
  ↓
k3d load balancer
  ↓
node:30080 (NodePort)
  ↓
Pod container:3000
\end{verbatim}

Portanto, o balanceador de carga do k3d permite o acesso ao host local
durante o desenvolvimento.

\subsubsection{Sobre o roteamento do
Kubernetes}\label{sobre-o-roteamento-do-kubernetes}

O cluster foi criado da seguinte forma:

\begin{Shaded}
\begin{Highlighting}[]
\ExtensionTok{k3d}\NormalTok{ cluster create }\AttributeTok{{-}{-}port}\NormalTok{ 8082:30080@agent:0 }\AttributeTok{{-}p}\NormalTok{ 8081:80@loadbalancer }\AttributeTok{{-}{-}agents}\NormalTok{ 2}
\end{Highlighting}
\end{Shaded}

Portanto, o \texttt{localhost:8082} é encaminhado para a porta 30080 do
agente de nó \texttt{:0} dentro do cluster. E se o aplicativo estiver
sendo executado em outro nó do cluster? O \texttt{@agent:0} da
configuração acima determina qual nó recebe o tráfego externo vindo da
sua máquina host. Depois disso, a rede do serviço Kubernetes assume o
controle e roteia automaticamente para o local correto do pod.

É por isso que os serviços NodePort são tão poderosos: você não precisa
saber ou se importar em qual nó o pod está. O serviço lida com todo o
roteamento internamente.

\begin{quote}
\textbf{Exercício 1.6 --- O projeto, etapa 4}

\textbf{Instruções:} Use um serviço NodePort para habilitar o acesso ao
projeto.
\end{quote}

\subsubsection{O que é um Ingress?}\label{o-que-uxe9-um-ingress}

Há ainda um recurso adicional que nos ajudará a servir o aplicativo: o
Ingress.

O recurso de acesso à rede de entrada,
\href{https://kubernetes.io/docs/concepts/services-networking/ingress/}{Ingress}
, é um tipo de recurso completamente diferente dos Serviços. Se você tem
o \href{https://en.wikipedia.org/wiki/OSI_model}{modelo OSI} memorizado,
o Ingress funciona na camada 7, enquanto os serviços funcionam na camada
4. Você pode vê-los usados juntos: primeiro o LoadBalancer mencionado
anteriormente e depois o Ingress para lidar com o roteamento. No nosso
caso, como não temos um balanceador de carga disponível, podemos usar o
Ingress como primeira parada. Se você está familiarizado com proxies
reversos como o Nginx, o Ingress deve parecer familiar.

Os Ingress são implementados por vários ``controladores'' diferentes.
Isso significa que os Ingress não funcionam automaticamente em um
cluster, mas lhe dão a liberdade de escolher qual controlador Ingress
funciona melhor para você. O K3s já tem o
\href{https://traefik.io/traefik}{Traefik} instalado. Outras opções
incluem Istio e o Nginx Ingress Controller, entre outras, mais
\href{https://kubernetes.io/docs/concepts/services-networking/ingress-controllers/}{aqui}.

Mudar para o Ingress exigirá que criemos um recurso do tipo Ingress. O
Ingress roteará o tráfego de entrada para um Serviço, mas o antigo
Serviço NodePort não fará isso, então o removeremos:

\begin{Shaded}
\begin{Highlighting}[]
\ExtensionTok{$}\NormalTok{ kubectl delete }\AttributeTok{{-}f}\NormalTok{ manifests/service.yaml}
\ExtensionTok{service} \StringTok{"hashresponse{-}svc"}\NormalTok{ deleted}
\end{Highlighting}
\end{Shaded}

Um recurso de Serviço do tipo
\href{https://kubernetes.io/docs/concepts/services-networking/service/\#type-clusterip}{ClusterIP}
fornece ao Serviço um IP interno que estará acessível dentro do cluster.

A seguinte definição de \texttt{service.yaml} permitirá o tráfego TCP da
porta 2345 para a porta 3000:

\begin{Shaded}
\begin{Highlighting}[]
\FunctionTok{apiVersion}\KeywordTok{:}\AttributeTok{ v1}
\FunctionTok{kind}\KeywordTok{:}\AttributeTok{ Service}
\FunctionTok{metadata}\KeywordTok{:}
\AttributeTok{  }\FunctionTok{name}\KeywordTok{:}\AttributeTok{ hashresponse{-}svc}
\FunctionTok{spec}\KeywordTok{:}
\AttributeTok{  }\FunctionTok{type}\KeywordTok{:}\AttributeTok{ ClusterIP}
\AttributeTok{  }\FunctionTok{selector}\KeywordTok{:}
\AttributeTok{    }\FunctionTok{app}\KeywordTok{:}\AttributeTok{ hashresponse}
\AttributeTok{  }\FunctionTok{ports}\KeywordTok{:}
\AttributeTok{    }\KeywordTok{{-}}\AttributeTok{ }\FunctionTok{port}\KeywordTok{:}\AttributeTok{ }\DecValTok{2345}
\AttributeTok{      }\FunctionTok{protocol}\KeywordTok{:}\AttributeTok{ TCP}
\AttributeTok{      }\FunctionTok{targetPort}\KeywordTok{:}\AttributeTok{ }\DecValTok{3000}
\end{Highlighting}
\end{Shaded}

O segundo recurso que precisamos é o novo Ingress, que será definido no
arquivo \texttt{ingress.yaml}. Ele precisa:

\begin{itemize}
\tightlist
\item
  Declarar que deve ser um Ingress
\item
  Encaminhar todo o tráfego para o nosso serviço
\end{itemize}

A configuração é assim:

\begin{Shaded}
\begin{Highlighting}[]
\FunctionTok{apiVersion}\KeywordTok{:}\AttributeTok{ networking.k8s.io/v1}
\FunctionTok{kind}\KeywordTok{:}\AttributeTok{ Ingress}
\FunctionTok{metadata}\KeywordTok{:}
\AttributeTok{  }\FunctionTok{name}\KeywordTok{:}\AttributeTok{ dwk{-}material{-}ingress}
\FunctionTok{spec}\KeywordTok{:}
\AttributeTok{  }\FunctionTok{rules}\KeywordTok{:}
\AttributeTok{  }\KeywordTok{{-}}\AttributeTok{ }\FunctionTok{http}\KeywordTok{:}
\AttributeTok{      }\FunctionTok{paths}\KeywordTok{:}
\AttributeTok{      }\KeywordTok{{-}}\AttributeTok{ }\FunctionTok{path}\KeywordTok{:}\AttributeTok{ /}
\AttributeTok{        }\FunctionTok{pathType}\KeywordTok{:}\AttributeTok{ Prefix}
\AttributeTok{        }\FunctionTok{backend}\KeywordTok{:}
\AttributeTok{          }\FunctionTok{service}\KeywordTok{:}
\AttributeTok{            }\FunctionTok{name}\KeywordTok{:}\AttributeTok{ hashresponse{-}svc}
\AttributeTok{            }\FunctionTok{port}\KeywordTok{:}
\AttributeTok{              }\FunctionTok{number}\KeywordTok{:}\AttributeTok{ }\DecValTok{2345}
\end{Highlighting}
\end{Shaded}

Então podemos aplicar tudo e visualizar o resultado:

\begin{Shaded}
\begin{Highlighting}[]
\ExtensionTok{$}\NormalTok{ kubectl apply }\AttributeTok{{-}f}\NormalTok{ manifests}
\ExtensionTok{ingress.networking.k8s.io/dwk{-}material{-}ingress}\NormalTok{ created}
\ExtensionTok{service/hashresponse{-}svc}\NormalTok{ configured}

\ExtensionTok{$}\NormalTok{ kubectl get svc,ing}
\ExtensionTok{NAME}\NormalTok{                       TYPE        CLUSTER{-}IP   EXTERNAL{-}IP   PORT}\ErrorTok{(}\ExtensionTok{S}\KeywordTok{)}    \ExtensionTok{AGE}
\ExtensionTok{service/kubernetes}\NormalTok{         ClusterIP   10.43.0.1    }\OperatorTok{\textless{}}\NormalTok{none}\OperatorTok{\textgreater{}}\NormalTok{        443/TCP    5m22s}
\ExtensionTok{service/hashresponse{-}svc}\NormalTok{   ClusterIP   10.43.0.61   }\OperatorTok{\textless{}}\NormalTok{none}\OperatorTok{\textgreater{}}\NormalTok{        2345/TCP   45s}

\ExtensionTok{NAME}\NormalTok{                                             CLASS    HOSTS   ADDRESS                            PORTS   AGE}
\ExtensionTok{ingress.networking.k8s.io/dwk{-}material{-}ingress}   \OperatorTok{\textless{}}\NormalTok{none}\OperatorTok{\textgreater{}}   \PreprocessorTok{*}\NormalTok{       172.21.0.3,172.21.0.4,172.21.0.5   80      16s}
\end{Highlighting}
\end{Shaded}

Podemos ver que o Ingress está escutando na porta 80. Como já abrimos
essa porta, podemos acessar o aplicativo em http://localhost:8081.

O fluxo de tráfego ao acessar o aplicativo agora é o seguinte:

\begin{verbatim}
http://localhost:8081/
  ↓
[Ingress Controller] (port 8081 mapped from host)
  ↓
[Ingress Rules] (path: / → hashresponse-svc:2345)
  ↓
[Service] (port 2345 → targetPort 3000)
  ↓
[Pod] (container listening on port 3000)
  ↓
[Application]
\end{verbatim}

\begin{quote}
\textbf{Exercício 1.7 --- Acesso externo com Ingress}

\textbf{Instruções:} O aplicativo ``Log output'' atualmente gera um
registro de data e hora e uma sequência aleatória (que ele cria na
inicialização) para os logs.

Adicione um endpoint para solicitar o status atual (timestamp e a string
aleatória) e um Ingress para que você possa acessá-lo com um navegador.

Você pode simplesmente armazenar a string aleatória na memória.
\end{quote}

\begin{quote}
\textbf{Exercício 1.8 --- O projeto, etapa 5}

\textbf{Instruções:} Mude para usar Ingress em vez de NodePort para
acessar o projeto. Você pode excluir o Ingress do aplicativo ``Log
output'' para que eles não interfiram neste exercício. Analisaremos mais
sobre caminhos e roteamento no próximo exercício e, nesse ponto, você
pode configurar o projeto para ser executado com o aplicativo ``Log
output'' lado a lado.
\end{quote}

\begin{quote}
\textbf{Exercício 1.9 --- Mais serviços}

\textbf{Instruções:} Desenvolva um segundo aplicativo que simplesmente
responda com ``pong 0'' a uma solicitação GET e aumente um contador (o
0) para que você possa ver quantas solicitações foram enviadas. O
contador deve estar na memória para que possa ser reiniciado em algum
momento. Crie uma nova implantação para ele e faça com que ele
compartilhe a entrada com o aplicativo ``Log output''. Rotas de
solicitações direcionadas a \texttt{/pingpong} para ele.

Em exercícios futuros, esta segunda aplicação será chamada de
``aplicação de pingue-pongue''. Ela será usada com o aplicativo ``Log
output''.

O aplicativo de pingue-pongue precisará escutar solicitações em
\texttt{/pingpong}, então talvez você precise fazer alterações em seu
código. Isso pode ser evitado configurando o Ingress para reescrever o
caminho, mas isso fica como um exercício opcional. Você pode conferir a
\href{https://kubernetes.io/docs/concepts/services-networking/ingress/\#the-ingress-resource}{documentação
de Ingress do Kubernetes}.
\end{quote}

\subsubsection{API de gateway}\label{api-de-gateway}

No Capítulo 4, encontraremos outra maneira de acessar nosso cluster: a
API Gateway, que é um conjunto evolutivo de recursos e configurações
projetados para gerenciar o tráfego de rede dentro dos clusters do
Kubernetes. O objetivo é fornecer uma maneira mais expressiva e
extensível de definir como o tráfego deve ser roteado e controlado, em
comparação com as soluções descritas nesta página.

\end{document}
